\section{The Measurement Framework}
\paragraph{Design.} Platform is a fixed facet: two to four named platforms are neither random nor exchangeable, so the generalizability study treats posts, events and coders as random and scopes every conclusion to the platforms studied. \todo{model specification}
\paragraph{Invariance and DIF.} Emotion components are sparse categorical indicators, so invariance is tested in categorical-indicator models with threshold constraints, judged by effect sizes rather than chi-square at social-media $N$; an anchor-free DIF estimator \citep{chen2023} runs as robustness; alignment \citep{muthen2018} accompanies the exact sequence. Invariance results are treated as diagnostic evidence, not a licence \citep{robitzsch2023}.
\paragraph{Lexicon-artifact control.} Every test is replicated with a non-lexicon classifier; only convergent non-invariance is interpreted as a mode effect \citep{ludwig2022}.
\paragraph{Multilevel structure.} Posts nest in events; invariance is examined in a two-level framework \citep{jak2013}, and the multilevel model includes time-since-event so platform reaction latency is part of the mode effect, not noise.
\blocked{RESULTS}
